\section{The unified taxonomy}\label{sec:taxonomy}

Harness design decomposes into 38 separately coded choices in nine layers, not into a small number of
architectures. Figure~\ref{fig:taxonomy-map} draws the structure and Table~\ref{tab:dimensions} lists
every dimension with the size of its value set. Each dimension is a question a reader can put to any
system, answer from a pinned artifact, and support with a cited line. The coded corpus fits this premise
better than a family structure: clustering falls short of both pre-stated floors for silhouette and
bootstrap stability (verdict negative), and on pairs where both dimensions carry a value the mean
corrected Cram\'er's $V$ is 0.167 across 666 pairs (\S\ref{sec:landscape}), which is what a space of
loosely related choices looks like when it is measured.

\begin{figure}[t]
\centering
\includegraphics[width=\linewidth]{taxonomy_map}
\caption{The unified taxonomy, generated from the frozen schema so it cannot drift from it. Top row: the
three necessary clauses, whose composition is the step loop
$\mathbf{A}\rightarrow\mathbf{B}\rightarrow\mathbf{C}\rightarrow\mathbf{A}$. Second row:
clauses~(iv)--(viii), which may be trivial; a trivial value is still a value, never
\texttt{not\_reported}. \textbf{M} carries metadata, not a clause. Dashed arrows: the deliberate
displacements of \textbf{F1} (clause~(iii)) and \textbf{G4} (clause~(vi)). Value counts as in
Table~\ref{tab:dimensions}.}
\label{fig:taxonomy-map}
\end{figure}

% Generated from the frozen schema by the Section 5 table generator; do not edit by hand.
% One float: main.tex does not load longtable, and the 38 single-line rows fit one float.
\begin{table}[tbp]
\caption{The 38 dimensions of the frozen schema (v1.0.0), grouped by layer. \emph{Values}: $n$ is
the size of a single-valued enumeration; $\{n\}$ is the size of a multi-valued one, whose cell holds a
set; a type name marks the four dimensions that are not enumerations. Full value sets, the judgment
call for each layer, and two worked examples are in Supplement~S1.}
\label{tab:dimensions}
\footnotesize
\begin{tabular}{@{}l l l l r@{}}
\toprule
Layer & ID & Key & What it records & Values \\
\midrule
A\enspace Context & A1 & \texttt{system\_prompt\_style} & how the per-call system prompt is produced & 4 \\
\phantom{A}\enspace assembly & A2 & \texttt{env\_context\_strategy} & how the model learns the repository or environment & 6 \\
 & A3 & \texttt{context\_compaction} & mechanisms that shrink the history sent & $\{6\}$ \\
 & A4 & \texttt{observation\_format} & format in which tool results reach the model & $\{5\}$ \\
\midrule
B\enspace Tool interface & B1 & \texttt{tool\_call\_format} & syntax by which the model requests actions & $\{6\}$ \\
 & B2 & \texttt{tool\_count} & tool definitions exposed at the default configuration & \emph{integer} \\
 & B3 & \texttt{edit\_primitive} & how the model expresses a file edit & $\{6\}$ \\
 & B4 & \texttt{tool\_schema\_source} & where the tool schema comes from & $\{4\}$ \\
 & B5 & \texttt{protocol\_standardization} & standard agent or tool protocols spoken & $\{4\}$ \\
\midrule
C\enspace Control loop & C1 & \texttt{loop\_primitives} & control-flow patterns driving calls and actions & $\{7\}$ \\
 & C2 & \texttt{planning\_granularity} & whether and how a plan is represented & 4 \\
 & C3 & \texttt{multi\_agent\_topology} & arrangement of agents at the default & 6 \\
 & C4 & \texttt{delegation\_mechanism} & how work is handed to another agent & 4 \\
 & C5 & \texttt{human\_in\_loop} & points where a human gates or steers & $\{5\}$ \\
\midrule
D\enspace Memory and & D1 & \texttt{short\_term\_state} & working state kept within a run & 3 \\
\phantom{D}\enspace state & D2 & \texttt{long\_term\_memory} & information persisted across runs and re-injected & $\{5\}$ \\
 & D3 & \texttt{state\_persistence} & whether an interrupted run can continue & 3 \\
\midrule
E\enspace Verification & E1 & \texttt{self\_verification} & checks the harness applies to the work & $\{6\}$ \\
\phantom{E}\enspace and repair & E2 & \texttt{retry\_policy} & task-level retry after a failed attempt & 4 \\
 & E3 & \texttt{rollback} & mechanism that reverts changes made in the run & 3 \\
\midrule
F\enspace Budget and & F1 & \texttt{termination\_condition} & conditions that end a run & $\{6\}$ \\
\phantom{F}\enspace termination & F2 & \texttt{cost\_controls} & mechanisms that limit or reduce spend & $\{4\}$ \\
 & F3 & \texttt{timeouts} & time limits the harness enforces & $\{3\}$ \\
\midrule
G\enspace Sandbox and & G1 & \texttt{execution\_isolation} & boundary within which actions execute & 5 \\
\phantom{G}\enspace environment & G2 & \texttt{filesystem\_access} & what the agent may touch inside G1 & 4 \\
 & G3 & \texttt{network\_policy} & network egress available to actions & 3 \\
 & G4 & \texttt{permission\_model} & how individual actions are authorized & 4 \\
\midrule
H\enspace Observability & H1 & \texttt{tracing} & richest run record the harness can emit & 4 \\
\phantom{H}\enspace and governance & H2 & \texttt{replayability} & whether a recorded trajectory can be re-run & 3 \\
 & H3 & \texttt{eval\_hooks} & benchmark evaluation wired into its own repositories & 2 \\
 & H4 & \texttt{guardrails} & safety controls on inputs, outputs, and actions & $\{4\}$ \\
\midrule
M\enspace Meta & M1 & \texttt{target\_domain} & task domains the harness is built for & $\{6\}$ \\
 & M2 & \texttt{open\_source} & whether the source is open & 3 \\
 & M3 & \texttt{model\_agnostic} & whether it runs across model providers & 2 \\
 & M4 & \texttt{primary\_artifact} & artifact the pinned version is coded from & 3 \\
 & M5 & \texttt{first\_release\_date} & first tagged release of the coded line & \emph{date} \\
 & M6 & \texttt{pinned\_version} & tag and commit per coded repository & \emph{string} \\
 & M7 & \texttt{stars} & flagship-repository stars, date-stamped in evidence & \emph{integer} \\
\bottomrule
\end{tabular}
\end{table}

Table~\ref{tab:dimensions} shows an instrument built mostly from small closed lists. Nineteen dimensions
are single-valued enumerations, fifteen hold a set, and four are an integer, a date, or a string; the
enumerations range from two to seven values. Every cell is therefore checkable against a closed list or a
stated counting, dating, or pinning rule, and the table is generated from the schema so that it cannot
drift from what was coded.

A layer is the set of dimensions that answers one clause of the working definition, and layers A to H
follow the clauses in order (\S\ref{subsec:layers}). Two placements depart from that map on purpose. The
termination condition sits in F so that termination policy and budget are reported together, and the
permission model sits in G because in practice it is configured alongside isolation. The ninth layer, M,
is metadata that stratification, weighting, and the time series need. Within a layer, one dimension records one
decision, made at one identifiable place in the pinned artifact, whose value a second reader can
recover from the same evidence. Per-dimension reliability is the test of that criterion: a dimension
that bundles two decisions yields readings that agree on the evidence and disagree on the value.

That criterion is why our cuts differ from prior taxonomies. Guo et al.\ place stopping criteria under
the control loop and budget control under verification and governance~\citep{guo2026qa2task}; a coder
cannot apply that split reliably, so we separate C from F, and E from G4 and H4. Li et al.'s context
layer covers both what the model sees this step and what persists between runs, which we split into A
and D~\citep{li2026harnessengineering}. Nine of Rombaut's twelve source-code dimensions map onto ours,
and that work supplied our evidence standard of a pinned commit with a file and
line~\citep{rombaut2026insidescaffold}. The other three have no dimension of ours to map to;
\S\ref{sec:related} names them and \S\ref{sec:threats} weighs what their omission costs.

\subsection{Multi-valued dimensions}\label{subsec:multivalued}

Fifteen dimensions hold a set, because on them a harness does several things at once: it stacks
compaction mechanisms, ships several call parsers, or gates a run at more than one point. A set cell
records every value reachable through shipped configuration, with the default named in the note. Where a value set contains \texttt{none}, it is a singleton by rule and never appears
beside another value. The registered reliability measure, exact-set agreement, is harsh on these cells:
two readings that agree on all but one primitive score as a complete disagreement, so the penalty grows
with the size of the true set rather than the size of the error. Loop primitives (C1) bears the cost. It
is the one dimension below the 0.6 freeze threshold on exact-set Cohen's kappa (0.590), and it clears
the threshold per value (0.731); Supplement~S1 reports all four measures, each of them
reproducibility between two model readings rather than correctness. We keep C1 and flag it as the
weakest dimension rather than redefine it, because redefining a dimension after seeing its reliability
is the post-hoc repair this paper audits others for.

\subsection{The absence rule}\label{subsec:absence}

The taxonomy separates a design decision from a documentation gap, and the absence rule enforces the
separation. An absence value, such as \texttt{none} or \texttt{open} for network policy, is a value on a
dimension: a choice the builders made. \texttt{not\_reported} is a property of the documents: the
sources were read and are silent. The rule (Table~\ref{tab:absence-rule}) codes an absence value only
when the evidence contains an enumerated place where the feature would be declared, and the feature is
not there; prose that never mentions a feature is never evidence that it is absent. A taxonomy that let
silence imply \texttt{none} would measure documentation while claiming to measure design, and every
distribution over its values would mix the two with no way to separate them. The rule is what lets the
under-reporting result speak about the literature rather than the systems; its history as a protocol
amendment and its measured effect on reliability are in \S\ref{sec:methods}.

\subsection{Boundary cases and known ambiguities}\label{subsec:ambiguities}

Coding two systems end to end, SWE-agent v1.1.0~\citep{yang2024sweagent} and OpenHands
v1.16.0~\citep{wang2024openhands}, raised 26 ambiguities, and each was resolved by a rule in the
coding manual rather than by changing a value set; Supplement~S1 lists all 26 with their resolutions.
Ten follow one of two patterns: code the shipped default and note the reachable alternatives (entries 1
and 13), or record the closest value and write the mismatch in the note (entries 3, 5, 6, 10, 12, 16, 24,
and 25). Entry 17 is the one departure from default-first coding: tracing is coded at the highest level
the artifact can reach, and the note says what runs by default. Table~\ref{tab:ambiguities} gives six.
Publishing the list is part of the instrument, because a reader who disagrees with a resolution can
recode that dimension from the released evidence strings.

% Six of the 26 manual entries; the full list is in the coding-manual supplement (S1).
\begin{table}[tbp]
\caption{Six of the 26 ambiguities recorded while coding the two worked systems, numbered as in the
manual: the two general rules (1, 26), the one departure from default-first coding (17), and three that
settle checklist dimensions (\S\ref{subsec:checklist}). All 26 are in Supplement~S1.}
\label{tab:ambiguities}
\footnotesize
\begin{tabular}{@{}r p{0.17\linewidth} p{0.31\linewidth} p{0.40\linewidth}@{}}
\toprule
\# & Dimension & Ambiguity & Resolution \tabularnewline
\midrule
1 & B2, C3--C5, D3, E1, E2, G1, G4 & \raggedright Capability shipped but off by default & \raggedright Code the default; name reachable alternatives in the note \tabularnewline
2 & B2 & \raggedright Toolsets, built-in, injected, and MCP tools & \raggedright One per schema sent; toolsets expanded; injected tools excluded \tabularnewline
10 & E3 & \raggedright Per-file tool-level undo & \raggedright \texttt{snapshot}; a repository reset is \texttt{git\_based} \tabularnewline
15 & G3 & \raggedright No network configuration anywhere & \raggedright \texttt{open} at medium confidence, citing the configuration surface \tabularnewline
17 & H1 & \raggedright Optional trace exporter versus always-on log & \raggedright Highest available level; default noted \tabularnewline
26 & General & \raggedright Evidence for an absence value & \raggedright The absence rule (\S\ref{subsec:absence}) \tabularnewline
\bottomrule
\end{tabular}
\end{table}

\subsection{Schema freeze}\label{subsec:freeze}

The schema is frozen at v1.0.0 with all 38 dimensions intact: none was added, removed, renamed, or
re-valued between the initial draft and the freeze on 2026-09-23. The freeze was gated on
per-dimension reliability rather than on a date, and the gate, Cohen's kappa of at least 0.6, was met on
37 of 38 dimensions. The one change to the cell contract came before the freeze and added
\texttt{unresolved} beside \texttt{not\_reported}, so that a cell the coder could not settle is never
pooled with a source that is silent. The released schema changelog records every change since the draft, and any
later change needs a changelog line and a registration update, so that the released dataset and this
section describe one instrument. Supplement~S1 carries the per-layer value sets, judgment calls,
decision rules, and worked examples.
